\section{材料与方法}

\subsection{模型冻结与可追溯性}

全文只评价原始检查点 \path{frankenstein_v28.pt}，SHA-256 为 \checkpointsha。检查点包含 ProteinMPNN 风格的共享结构条件编码/解码主干、20 类天然氨基酸输出头，以及按天然残基类别选择的甲基化专家头。ProteinMPNN 的基本思想与原始实现见 Dauparas 等\cite{dauparas2022proteinmpnn}；本研究所述 N-甲基化扩展与其结果均以项目仓库中的冻结代码、表格和哈希为准。

检查点顶层只保存 \path{model_state_dict}，没有 epoch、随机种子、优化器状态、训练数据 manifest 或输入权重哈希。合并脚本的默认参数指向一个 V28 基础模型和一个甲基化专家模型，但没有运行 manifest 能证明最终文件正是由这两个默认路径生成。因此，本文可以冻结并评价最终 checkpoint，却不虚构无法追溯的完整训练谱系，也不宣称相对初始化模型的增益。

本文未为写作重新运行模型、HighFold/ESMFold、TM-align、通透性预测或 PyRosetta。所有数值均从此前已经生成并审计的位点预测表、FASTA、结构匹配表和 RMSD 表中读取。此策略减少了因软件环境、随机种子或上游服务变化而引入的第二套结果。源文件与哈希见 \path{SOURCE_MANIFEST.json}。

\subsection{单体数据与来源审计}

原始历史仓库在冻结提交中保留了 751 条合并单体记录，并以 600/151 划分训练集与测试集。对对应 751 个原始 PDB 的逐文件复核显示：751/751 的 \texttt{EXPDTA} 字段均标为 \texttt{THEORETICAL MODEL}，并且都含有 \texttt{MODEL GENERATED BY ROSETTA} 记录。其中 406 条 \path{Me_} 前缀来自 Rosetta \path{2023.35+release.23439d3}，345 条 \path{pdb_} 前缀来自 Rosetta \path{2025.25+release.a0cefad01b}。因此，这批单体应准确称为两个版本的 Rosetta 理论结构面板，而不是 Baker 实验真结构。现有精确名称与天然化序列检查没有发现 train--test 完全相同记录，但尚未做同源聚类，因此不能排除较远的同源泄漏。

\begin{table}[H]
\centering
\caption{单体数据文件的可核验事实。所有 751 个原始 PDB 均标记为 Rosetta 生成的理论模型；规划截图中的“Baker 33、Rosetta 400+300”与冻结文件不一致。}
\label{tab:monomer-data}
\begin{tabularx}{\textwidth}{L{31mm}C{16mm}C{16mm}C{16mm}L{44mm}Y}
\toprule
数据对象 & 总条数 & \path{Me_} & \path{pdb_} & Git blob & 本文用途 \\
\midrule
历史合并集 & 751 & 406 & 345 & \path{768d7863ac1a...} & Rosetta-2023/2025 理论结构 \\
历史训练集 & 600 & 323 & 277 & \path{720b7e65d4bf...} & 训练来源背景；无完整训练日志 \\
留出测试集 & 151 & 83 & 68 & \path{bd18c23881cb...} & Rosetta-2023/2025 为83/68条，共1505位点 \\
Ser修正测试集 & 151 & 83 & 68 & SHA-256 \path{56f877bb9987...} & 论文主真值 \\
\bottomrule
\end{tabularx}
\end{table}

旧测试标签包含 323 个小写甲基化位点，其中 74 个为 Ser。后续来源审计确认 62 个 Ser 应为天然 Ser；修正后为 261 个甲基化阳性、1244 个阴性，甲基化率为 17.34\%。模型权重、每个位点预测概率与天然氨基酸目标均未改变，只有评价真值改变。因此，本稿主结果使用修正标签；旧 323 阳性口径仅用于敏感性审计。

751个 PDB 的文件级来源、Rosetta版本、训练/测试归属、SHA-256 与 Git blob 逐行列于 \path{data/Monomer_PDB_Provenance_751.csv}；扫描程序为 \path{scripts/audit_monomer_pdb_provenance.py}。精确名称与天然化序列的 train--test 重叠审计结果位于 \path{data/Monomer_Train_Test_Exact_Overlap_Audit.json}，两类交集均为0；该检查不等价于同源聚类，不能排除较远同源泄漏。

\subsection{单体序列与甲基化指标}

天然化氨基酸恢复率定义为
\begin{equation}
\mathrm{RAA}=\frac{1}{N}\sum_{i=1}^{N}\mathbb{1}\!\left(\widehat a_i=a_i\right),
\end{equation}
其中所有小写 N-甲基化 token 在比较前映射回对应大写天然氨基酸。

甲基化二分类采用两种输入条件：
\begin{description}
  \item[已知天然序列（known-sequence）] 用真实天然氨基酸类别选择相应专家头，隔离甲基化头自身的辨别能力；
  \item[端到端（end-to-end）] 先由基础头预测天然氨基酸，再用预测类别选择专家头，包含基础序列错误的级联影响。
\end{description}

二者均在 1505 个真实位点上计算 ROC-AUC、Accuracy、Precision、Recall、F1 与假阳性率。预设甲基化判定为 $p_i>0.6$，等号不通过。尚哥清单中的 ``RMC'' 不是统一的标准缩写；若它指不考虑天然氨基酸类别是否正确的甲基阳性二分类恢复率，则其数学定义就是 Recall：
\begin{equation}
\mathrm{Recall}=\frac{TP}{TP+FN}.
\end{equation}
若 ``RMC'' 指“天然氨基酸身份和甲基状态同时正确的真实甲基残基比例”，则应另记为 exact methylated-residue recovery；若以全部位点为分母，则是联合扩展 token recovery。本文并列给出三者的清楚名称和分母，不把它们混成一个 RMC 数字。

\subsection{17 个天然复合物与全温度生成池}

天然复合物集合含17个靶点，共251个选中短肽位点。天然序列评价保留每个结构中预先选定的全部短肽链，因此部分靶点贡献多条等价肽链；这与生成结构中每个候选只评价一条预测末链的口径不同。该天然集合没有甲基化阳性标签，因此 AUC、Precision、Recall 与 F1 在数学上不可定义；只能报告基础 RAA、阴性位点上的假阳性率与预测阳性率。

原始生成池覆盖温度 0.01、0.1、0.2、0.3 与 0.5。记录总数、去重数、天然氨基酸恢复和甲基标记率均在每个温度内汇总。此前的 best85 是对每个靶点、每个温度各选一条，共 $17\times5=85$ 条；选择标准使用了天然序列恢复率，因此属于 \textbf{oracle best-of-N}，不能作为未知天然序列时的无偏部署性能。

\subsection{六个非 3AV 复合物的预设主分析}

主分析只包括 1SFI、3P8F、3WNE、3ZGC、4K1E 和 4KEL，排除全部 3AV 系列。候选只来自主目录温度 $T=0.5$ 的 FASTA；排除 \path{.ipynb_checkpoints} 等重复/缓存目录。原始序列共 544 条。甲基化保留门要求序列至少含一个由原始生成器以严格 $p>0.6$ 保存的小写标记，得到 476 条。

“有效生成概率”用两个分母同时报告：
\begin{align}
\text{保留后条件通过率} &= \frac{N_{\mathrm{RMSD\ pass}}}{476},\\
\text{原始端到端产率} &= \frac{N_{\mathrm{RMSD\ pass}}}{544}.
\end{align}
前者描述进入结构阶段后的质量，后者回答一次原始生成得到合格候选的概率；二者不可互换。

\subsection{六复合物 RMSD 的冻结定义}

对每个预测复合物与相应天然复合物仅进行一次全复合物 CA 对齐。随后在同一个变换后的坐标系中计算：

\begin{enumerate}
  \item \textbf{全局复合物 CA RMSD：}使用已审计的复合物匹配 CA 原子；
  \item \textbf{环肽 CA RMSD：}使用预测和天然结构的完整末链肽 CA，枚举所有正向循环移位并取最小值；
  \item 不允许反向序列，不对肽进行第二次自对齐，不丢弃难配对位点；完整位点覆盖是硬质量门。
\end{enumerate}

联合 $<d$ 判定要求全局复合物与环肽 RMSD 均严格小于 $d$，其中 $d\in\{3,5\}$~\AA。循环肽没有唯一的线性起点，因此 best-forward 循环移位是主定义；固定起点结果只作为表示敏感性分析。

\subsection{best85 的探索性结构与下游指标}

best85 的旧结构分析与六复合物主分析使用不同队列和 RMSD 定义，单独报告：

\begin{itemize}
  \item \textbf{受体框架对齐后肽 RMSD：}先拟合受体 CA，再测肽 CA 或 N/CA/C backbone，反映结合姿态；
  \item \textbf{肽自对齐 RMSD：}单独拟合肽本身，反映内部形状，不反映结合位置；
  \item \textbf{多样性：}同一靶点五个温度代表之间两两比较，共170对，使用对称 TM-score 的 $1-\mathrm{TM}$；TM-score 的尺寸归一化思想见 Zhang 与 Skolnick\cite{zhang2004tmscore}；
  \item \textbf{置信度：}优先区分 PDB CA B-factor 均值代理与 HighFold COMMENT 中的 global/chain pLDDT、pTM、ipTM、PAE 字段。结构预测模型的置信度概念可参照 AlphaFold 与 ESMFold 文献\cite{jumper2021alphafold,lin2023esmfold}，但本项目字段不应跨实现直接等同；
  \item \textbf{通透性：}只报告外部预测值及其分布，无实验标签；
  \item \textbf{能量：}将非标准甲基残基自然化后，在固定构象上计算 Rosetta 分数。没有 FastRelax、最小化、repacking 或显式 N-甲基参数。单位为 Rosetta Energy Unit（REU），不是 kcal/mol；cross-interface 分数不是结合自由能\cite{chaudhury2010pyrosetta,leaverfay2011rosetta3}。
\end{itemize}

\subsection{统计与缺失值规则}

比例同时给出计数与分母；关键联合 RMSD 比例给出 Wilson 95\% 置信区间。偏态指标优先报告中位数与四分位距。所有阈值采用严格不等号。数学上不可定义、未计算或源字段缺失的指标标记为“不可报告/未估计”，绝不以 0 代替。没有原始成对单体结构结果，因此单体 RMSD、pLDDT、能量和通透性均不作数值报告。
